%\begin{center}
%\large \bf \runtitulo
%\end{center}
%\vspace{1cm}
\chapter*{\runtitulo}

\noindent El \emph{morphing tímbrico}, la generación de timbres intermedios
entre sonidos conocidos, es una herramienta de interés para la producción
musical y el diseño sonoro, ya que permite expandir la paleta sonora más allá
de los sonidos tradicionales. Las técnicas usuales operan en el
espacio latente de un único modelo generativo entrenado sobre múltiples
timbres; si bien un modelo de gran capacidad puede representar varios de
ellos con fidelidad, en modelos pequeños esta estrategia compromete la
fidelidad de cada uno.

En este trabajo se propone un enfoque alternativo: entrenar autoencoders
variacionales ($\beta$-VAE) especializados por timbre a partir de un
modelo base común, y generar timbres híbridos mediante la interpolación
de los pesos de esos modelos.

El sistema se entrena sobre espectrogramas de magnitud obtenidos por STFT
de cuatro fuentes sonoras (piano, guitarra, voz y bajo) y se evalúa mediante
Fréchet Audio Distance (FAD), similitud coseno sobre \emph{embeddings} MERT,
análisis del espacio latente con UMAP y error de reconstrucción.

Los resultados muestran que especializar cada modelo mediante \emph{fine-tuning}
a partir de un modelo base común —en lugar de entrenarlo desde cero— produce
transiciones tímbricas más graduales y predecibles, en especial al combinarse
con la regularización del espacio latente ($\beta$-VAE). Al apoyarse en modelos
pequeños y livianos, el enfoque resulta además potencialmente apto para
integrarse en un sistema de síntesis en tiempo real.

\bigskip

\noindent\textbf{Palabras claves:} morphing tímbrico, autoencoder variacional,
$\beta$-VAE, interpolación de pesos, síntesis de audio, Fréchet Audio
Distance.
